\section{The main result: Proposition 2}
\label{sec:06_main_result}

Proposition~2 compares three economies with the same classes, costs, prior, and order slot.
Only the incumbent distribution changes.
Entry equals its floor against the weak incumbent.
It exceeds its floor against the strong incumbent.
It returns to its floor against the very strong incumbent.

\subsection{The statement and the role of each hypothesis}
\label{sec:06_main_result-statement}

Let $F_1$ strengthen $F_0$ by first-order stochastic dominance, and let $F_2$ strengthen $F_1$.
The distributions are continuous.
The classes satisfy the class ranking: $G_H$ dominates $G_L$, and every class value meets the reserve.
Write
\[
B_i=B_{F_i},\quad \Delta_i=\DT(F_i),\quad M=M(\pi),\quad m=m(\pi),\quad
\kappa=\min\{m,1-M\}.
\]
Keep $0<\pi,\rho,\lambda<1$, $k>0$, and the slot restriction $0<\beta\le1/3$.

\begin{resultbox}[Proposition 2: competition creates competition]
Assume the following three strict hypotheses:
\begin{align*}
\text{(H1)}&\qquad B_0(\pi)<c_H<B_1(M),\\
\text{(H2)}&\qquad \Delta_0<k<\rho\kappa\Delta_1,\\
\text{(H3)}&\qquad B_2(M)<c_H.
\end{align*}
Each economy has a unique equilibrium outcome under truthful bids. \paper{Prop.~2}
\begin{enumerate}[label=(\roman*),leftmargin=2em]
\item At $F_0$, both classes choose zero.
Every posterior is $\pi$, every price is the same, and entry is $\rho$.
\item At $F_1$, class $H$ buys and class $L$ sells with probability one.
The posteriors are $(m,\pi,M)$, and only a purchase recruits the high-cost challenger.
Entry is $\Ent_1=\rho+(1-\rho)(b+\pi\lambda)>\rho$.
High-class ownership rises from $\pi\rho W_H(F_0)$ to $\pi[\rho+(1-\rho)(b+\lambda)]W_H(F_1)$ if and only if
$[\rho+(1-\rho)(b+\lambda)]W_H(F_1)>\rho W_H(F_0)$.
\item At $F_2$, the orders remain full and the posteriors remain $(m,\pi,M)$.
The three prices differ, but entry after every order equals $\rho$.
\end{enumerate}
\end{resultbox}

The high-cost window (H1) sets the two preparation incentives.
At the weak prior, gross profit falls below the high cost.
At the strong top posterior, gross profit exceeds it.
The trading-cost window (H2) fixes orders against every candidate.
The ceiling (H3) stops high-cost preparation in the third economy.

The trading-cost window exists exactly when
\[
\Delta_0<\rho\kappa\Delta_1
\iff \frac{\Delta_0}{\Delta_1}<\rho\kappa.
\]
The ratio shows how much the target-payoff spread must rise.
Parts (i) and (ii) need only (H1) and (H2).
The proof below also works for every $\beta<1/2$.

\subsection{Step 0: derive the inequalities that the proof needs}
\label{sec:06_main_result-implied}

Proposition~1 gives $B_1(\pi)\le B_0(\pi)$.
Thus (H1) implies
\[
B_1(\pi)<c_H<B_1(M).
\]
Subtract the two gross profits:
\[
B_1(M)-B_1(\pi)=(M-\pi)(g_{H,1}-g_{L,1})>0.
\]
Since $M>\pi$, we obtain $g_{H,1}>g_{L,1}$.
Therefore $\tau_1$ is defined, and
\[
\pi<\tau_1=\frac{c_H-g_{L,1}}{g_{H,1}-g_{L,1}}<M.
\]
Also $m<\pi$, so neither a sale nor zero can recruit at $F_1$ once full orders are established.

Hypothesis (H2) gives $\Delta_1>k/(\rho\kappa)>0$.
Proposition~1 gives $\Delta_2\ge\Delta_1$, so
\[
\rho\kappa\Delta_2\ge\rho\kappa\Delta_1>k.
\]
The class ranking also gives $g_{H,2}\ge g_{L,2}$.
Hence $B_2(\mu)\le B_2(M)<c_H$ whenever $\mu\le M$.
The zero-cost challenger gives $e(x)\ge\rho$ in every candidate of all three economies.

\subsection{Step 1: force and verify the weak tuple}
\label{sec:06_main_result-weak}

Fix any candidate at $F_0$.
By \Cref{sec:05_investor-force}, $\Delta_0<k$ makes zero the unique best response in each class.
Thus every equilibrium has $\sigma_H=\sigma_L=\delta_0$, where $\delta_q$ is a point mass at $q$.

These orders give $f_H=f_L$ at every order.
Every order posterior, and every price belief, equals $\pi$.
Since $B_0(\pi)<c_H$, high-cost preparation is zero.
The only possible tuple is
\[
\mathcal E_0=(\delta_0,\delta_0,\mu_0,P_0,\bar\mu_0,\zeta_0),
\]
with truthful bids and
\[
\mu_0(x)=\pi,\quad e_0(x)=\rho,\quad
P_0(x)=\bar P_0:=t_{0,0}+\rho w_0(\pi),\quad
\bar\mu_0(\bar P_0)=\pi,\quad \zeta_0(\bar P_0)=0.
\]
Here $w_i(\mu)=t_{L,i}-t_{0,i}+\Delta_i\mu$.
Now check the six conditions against this fixed tuple. \paper{Definition~A.1; Appendix~A.3}
\begin{enumerate}[label=(E\arabic*),leftmargin=3em]
\item Since $f_H=f_L$, Bayes' rule gives $\mu_0(x)=\pi$.
\item The single price pools all orders.
Both conditional order laws sum to one, so its posterior is $\pi$.
\item The strict inequality $B_0(\pi)<c_H$ makes $\zeta_0=0$ optimal.
The zero-cost challenger prepares by the model's convention.
\item Conditional terminal value is $t_{0,0}+\rho(t_{j,0}-t_{0,0})$ in class $j$.
Its average at $\pi$ is exactly $\bar P_0$.
\item At this tuple, $K_H=\rho\Delta_0(1-\pi)<k$ and $K_L=\rho\Delta_0\pi<k$.
Thus zero uniquely beats both alternatives in either class.
\item Truthful bids are weakly dominant in the cash second-price auction with its reserve.
\end{enumerate}
This proves existence after the earlier bounds prove uniqueness.
Entry is $\rho$, and high-class ownership is $\pi\rho W_H(F_0)$.

\subsection{Steps 2 and 3: force and verify the strong tuple}
\label{sec:06_main_result-strong}

Fix any candidate at $F_1$.
By \Cref{sec:05_investor-force}, $k<\rho\kappa\Delta_1$ makes the correct order the unique best response in each class.
Hence every equilibrium has $\sigma_H=\delta_1$ and $\sigma_L=\delta_{-1}$.
The order laws now force the posteriors $(m,\pi,M)$.
Price sufficiency (\Cref{sec:04_slot_beliefs-sufficiency}) and the inequalities of Step~0 force entry $(\rho,\rho,1)$, with no tie.

The unique possible tuple is
\[
\mathcal E_1=(\delta_1,\delta_{-1},\mu_1,P_1,\bar\mu_1,\zeta_1),
\]
with truthful bids and, on $x=(-1,0,1)$,
\begin{gather*}
\mu_1=(m,\pi,M),\qquad e_1=(\rho,\rho,1),\\
P_1(x)=t_{0,1}+e_1(x)w_1(\mu_1(x)),\qquad
\bar\mu_1(P_1(x))=\mu_1(x),\qquad \zeta_1(P_1(x))=\1\{x=1\}.
\end{gather*}
The verification must establish the price beliefs directly.
It must not assume entry optimality to prove that same condition.
\begin{enumerate}[label=(E\arabic*),leftmargin=3em]
\item The likelihoods are $f_H=(b,b_0,b+\lambda)$ and $f_L=(b+\lambda,b_0,b)$.
Bayes' rule gives $(m,\pi,M)$.
\item The stated entry is nondecreasing in the posterior and is at least $\rho$.
The price-difference algebra in \Cref{sec:04_slot_beliefs-sufficiency} gives $P_1(-1)<P_1(0)<P_1(1)$ directly.
Each price identifies one order, so its belief equals that order's posterior.
\item We have $B_1(m)\le B_1(\pi)<c_H<B_1(M)$.
Thus precisely a purchase recruits the high-cost challenger.
\item At each order, average $t_{0,1}+e_1(x)(t_{j,1}-t_{0,1})$ with its Bayes weights.
This equals the stated $P_1(x)$.
\item Direct substitution gives $K_H=\Delta_1(1-M)>k$ and $K_L=\rho\Delta_1m>k$.
The candidate-wide lower bound implies both strict inequalities.
The unique best responses are therefore the stated full orders.
\item Truthful bids remain weakly dominant.
\end{enumerate}

\subsection{Step 4: entry, ownership, and the benchmark}
\label{sec:06_main_result-entry}

Only a purchase recruits a high-cost challenger.
In class $H$, a purchase occurs through informed presence or a liquidity purchase:
\[
\Pr(X=1\mid H)=\lambda+(1-\lambda)\beta=b+\lambda.
\]
In class $L$, only the liquidity purchase reaches it, so $\Pr(X=1\mid L)=b$.
Therefore
\begin{align*}
e_H&=\rho+(1-\rho)(b+\lambda),\\
e_L&=\rho+(1-\rho)b,\\
\Ent_1=\pi e_H+(1-\pi)e_L
&=\rho+(1-\rho)[\pi(b+\lambda)+(1-\pi)b]\\
&=\rho+(1-\rho)(b+\pi\lambda)>\rho.
\end{align*}
The difference $e_H-e_L=(1-\rho)\lambda>0$ shows the class selection.
Prices recruit both classes, but recruit class $H$ more often.

Let $W_H(F)=\E_H[F(\theta)]$ denote the probability that a prepared class-$H$ challenger wins.
The class supports lie above the reserve.
Preparation is independent of the realized values conditional on the class.
Thus the exact ownership comparison is
\begin{align*}
\Own_{H,1}>\Own_{H,0}
&\iff \pi e_HW_H(F_1)>\pi\rho W_H(F_0)\\
&\iff [\rho+(1-\rho)(b+\lambda)]W_H(F_1)>\rho W_H(F_0).
\end{align*}
More entry alone does not guarantee more high-class ownership under general classes.
A stronger incumbent can reduce $W_H$.

A useful sufficient condition is
\[
(1+\rho)W_H(F_1)\ge2\rho W_H(F_0)>0,
\qquad b+\lambda>\frac12.
\]
Indeed, $e_H>\rho+(1-\rho)/2=(1+\rho)/2$.
Multiplication by $W_H(F_1)>0$ gives
\[
e_HW_H(F_1)>\frac{1+\rho}{2}W_H(F_1)\ge\rho W_H(F_0).
\]
In the point-mass benchmark, $h>r_i$ gives $W_H(F_i)=1$.
The exact condition then reduces to $e_H>\rho$, which always holds here.

At $r_1=\reg{prim_r1}$, \Cref{tab:03_auction-table1} gives $g_L=1/8$ and $g_H=203/24$, so
$\tau_1=(6-1/8)/(25/3)=141/200$.
Since $\pi<\tau_1<M$, only the purchase recruits.
With $\beta=\reg{beta_B}$ and $u=\eu-1$, the script gives
\begin{align*}
\lambda&=\reg{lambda_B},& b&=\reg{b_B},& b+\lambda&=\reg{b_plus_lambda_B},\\
e_H&=\reg{eH_r1_B},& e_L&=\reg{eL_r1_B},& \Ent_1&=\reg{entry_r1_B}.
\end{align*}
Finally $\Own_{H,1}=\pi e_H=\reg{own_r1_B}$.
The weak values are $\Ent_0=\reg{entry_r0}$ and $\Own_{H,0}=\reg{own_r0}$.

\subsection{Step 5: the very strong tuple and the entry ceiling}
\label{sec:06_main_result-ceiling}

At $F_2$, $\Delta_2\ge\Delta_1$ preserves the strict incentive for both full orders under every candidate.
Hence the posteriors are again $(m,\pi,M)$.
Every price belief averages these three posteriors and is at most $M$.
Hypothesis (H3) therefore makes high-cost preparation strictly unprofitable after every price.
The only possible tuple is $\mathcal E_2=(\delta_1,\delta_{-1},\mu_2,P_2,\bar\mu_2,\zeta_2)$ with truthful bids,
\[
\mu_2=(m,\pi,M),\quad e_2\equiv\rho,\quad
P_2(x)=t_{0,2}+\rho w_2(\mu_2(x)),\quad
\bar\mu_2(P_2(x))=\mu_2(x),\quad \zeta_2\equiv0.
\]
The six checks repeat those of $\mathcal E_1$ with $e\equiv\rho$.
(E1) holds by the full-order likelihoods.
(E2): two prices differ by $\rho\Delta_2[\mu_2(x)-\mu_2(x')]\neq0$, so each price reveals its posterior.
(E3): $B_2(\mu_2(x))\le B_2(M)<c_H$ at every price.
(E4): the conditional terminal values average to $P_2(x)$.
(E5): $K_H=\rho\Delta_2(1-M)>k$ and $K_L=\rho\Delta_2m>k$.
(E6): truthful bids remain weakly dominant.

At $r_2=\reg{prim_r2}$, gross profit at the top posterior is
\[
B_{r_2}(M)=\frac5{48}+\frac{1451}{180}M
=\reg{p2_base_r2_low_profit_plus}<6.
\]
The cost margin is $\reg{margin_return_B}$.
Entry is $\rho$, yet the three prices still differ (\Cref{tab:07-outcomes}).
Information survives even though high-cost preparation stops.

For the uniform benchmark, expand gross profit at the top posterior:
\begin{align*}
B_r(M)
&=(1-M)\frac{\ell^2-p^2}{2r}
  +M\left(h-\frac r2-\frac{p^2}{2r}\right)\\
&=Mh-\frac{Mr}{2}+\frac{(1-M)\ell^2-p^2}{2r}.
\end{align*}
Each class profit falls strictly with $r$, so this weighted average also falls strictly.
The entry ceiling $r_C$ solves $B_{r_C}(M)=c_H$.
Multiplication by $2r_C$ gives
\[
Mr_C^2-2(Mh-c_H)r_C-[(1-M)\ell^2-p^2]=0.
\]
In the benchmark, $(1-M)\ell^2-p^2=1/(1+\eu)-1/4>0$, so the product of the two roots, $-[(1-M)\ell^2-p^2]/M$, is negative.
The positive root is
\[
r_C=\frac{Mh-c_H+\sqrt{(Mh-c_H)^2+M[(1-M)\ell^2-p^2]}}{M}
=\reg{rC_B}.
\]
Above $r_C$, even the top posterior fails to justify the high cost.
At $r_C$, a purchase leaves the high-cost challenger indifferent.
Rule C1 gives total entry $\Ent_1$ there; rule CM permits any value in $[\rho,\Ent_1]$.

\subsection{Derive the exact admissible band}
\label{sec:06_main_result-band}

Fix the benchmark auction, costs, and strengths, and let $\pi=1/2$.
The threshold $\tau(r)=(48r-3)/[4(20r-r^2-1)]$ of \Cref{sec:03_auction-derivatives} gives
\[
\tau(r_0)=\frac{195}{308},\qquad
\tau(r_1)=\frac{141}{200},\qquad
\tau(r_2)=\frac{4245}{5804},
\]
and \Cref{tab:03_auction-table1} gives $\Delta_0=1/60$, $\Delta_1=2/3$, and $\Delta_2=169/180$.
The weak-economy margins do not depend on the slot:
\[
c_H-B_0(1/2)=6-\frac{1153}{240}=\frac{287}{240}>0,\qquad
k-\Delta_0=\frac1{50}-\frac1{60}=\frac1{300}>0.
\]

At this prior, $M=(1+u)/(2+u)$ and $\kappa=1/(2+u)$.
The right side of (H1) becomes
\begin{align*}
M>\frac{141}{200}
&\iff 200(1+u)>141(2+u)\\
&\iff 59u>82
\iff u>\frac{82}{59}.
\end{align*}
The ceiling (H3) becomes
\begin{align*}
M<\frac{4245}{5804}
&\iff5804(1+u)<4245(2+u)\\
&\iff1559u<2686
\iff u<\frac{2686}{1559}.
\end{align*}
The right side of (H2) becomes
\begin{align*}
\frac1{50}<\frac1{6(2+u)}
&\iff6(2+u)<50
\iff u<\frac{19}{3}.
\end{align*}
Equivalently, it requires $M<22/25$.
Since $2686/1559<19/3$, the trading bound is redundant within the narrower band.
We therefore obtain the exact equivalence
\begin{resultbox}[The admissible band at the prior one half]
\[
\text{(H1)--(H3)}
\iff \frac{82}{59}<u<\frac{2686}{1559}
\iff \frac{141}{200}<M<\frac{4245}{5804}.
\]
This is Proposition~A.1(a). \paper{Prop.~A.1}
\end{resultbox}

At fixed $\beta$, $\lambda=u\beta/(1+u\beta)$ has derivative $\beta/(1+u\beta)^2>0$.
Thus the equivalent presence interval is
\[
\frac{82\beta}{59+82\beta}<\lambda<\frac{2686\beta}{1559+2686\beta}.
\]
At $\beta=\reg{beta_B}$, its endpoints are $\reg{lambda_lo_B}$ and $\reg{lambda_hi_B}$.
The benchmark $\lambda=\reg{lambda_B}$ lies strictly between them.

\subsection{The prior window and the boundary cases}
\label{sec:06_main_result-priors}

Now fix $1+u=\eu$ and vary the prior.
For any $0<\tau<1$, convert the top-posterior condition into prior odds:
\begin{align*}
M(\pi)>\tau
&\iff \frac{\pi}{1-\pi}\eu>\frac{\tau}{1-\tau}\\
&\iff \pi\eu(1-\tau)>\tau(1-\pi)\\
&\iff \pi>\frac{\tau}{\tau+(1-\tau)\eu}.
\end{align*}
The right side of (H1) and (H3) therefore give
\[
\frac{141}{141+59\eu}<\pi<\frac{4245}{4245+1559\eu}.
\]
The endpoints are approximately $\reg{prior_lo_B}$ and $\reg{prior_hi_B}$.
We must also check the other hypotheses.
The left side of (H1) requires $\pi<\tau(r_0)=195/308$.
The right side of (H2) requires
\[
\kappa(\pi)>\frac{k}{\rho\Delta_1}=\frac{1/50}{(1/4)(2/3)}=\frac3{25}.
\]
Both terms in the minimum must exceed $3/25$:
\begin{align*}
m(\pi)>\frac3{25}
&\iff 25\pi>3[\pi+(1-\pi)\eu]\\
&\iff\pi>\frac{3\eu}{22+3\eu},\\
1-M(\pi)>\frac3{25}
&\iff25(1-\pi)>3[1-\pi+\pi\eu]\\
&\iff\pi<\frac{22}{22+3\eu}.
\end{align*}
This interval has endpoints $\reg{prior_kappa_lo_B}$ and $\reg{prior_kappa_hi_B}$.
It contains both the full prior window and the window for parts (i) and (ii).
The full intersection therefore equals the first prior interval above.
Without (H3), parts (i) and (ii) require exactly
\[
\frac{141}{141+59\eu}<\pi<\frac{195}{308}.
\]

The slot band describes strict hypotheses, rather than every boundary outcome.
For $u<82/59$, (H2) still forces full orders at $r_1$, but $M<\tau(r_1)$.
Entry at $r_1$ is then $\rho$.
For $2686/1559\le u<19/3$, full orders persist at both stronger strengths.
A purchase at $r_2$ recruits under C1, so the return to the floor fails.
For $u\ge19/3$, the right side of (H2) fails; this fact alone specifies no complete equilibrium outcome.

At the lower boundary $u=82/59$, the posteriors satisfy $M=141/200$ and $m=59/200$.
Orders at $r_0$ remain uniquely zero because $\Delta_0-k=-1/300$.
Orders at $r_1$ and $r_2$ remain uniquely full because
\[
\rho\kappa\Delta_1-k
=\frac14\frac{59}{200}\frac23-\frac1{50}=\frac7{240}>0.
\]
Hypothesis (H3) remains strict.
A purchase at $r_1$ alone now creates a challenger tie.
C1 gives total entry $\Ent_1$ at that boundary slot.
CM allows $e(1)\in[\rho,1]$, so total entry is
\[
\Ent=\rho+[e(1)-\rho](b+\pi\lambda)\in[\rho,\Ent_1].
\]
At $\beta=\reg{beta_B}$, this boundary value of $\Ent_1$ is approximately $0.515737$. % verify/06_main_result.py
It differs from the benchmark value because the boundary slot has a different $u$ and $\lambda$.

\subsection{Why H1 needs the weak prior, and why the region is open}
\label{sec:06_main_result-scope}

The weaker condition $B_1(\pi)<c_H$ cannot replace $B_0(\pi)<c_H$.
Keep the benchmark slot and other primitives, but set $c_H=\reg{nec_cH}$ and $r_2=\reg{nec_r2}$.
The proposed weaker list has positive margins:
\begin{align*}
c_H-B_1(1/2)&=\frac{49}{120}>0,& B_1(M)-c_H&=\reg{nec_margin_recruit_B}>0,\\
k-\Delta_0&=\frac1{300}>0,& \rho\kappa\Delta_1-k&=\reg{nec_margin_kstrong_B}>0,\\
c_H-B_2(M)&=\reg{nec_margin_return_B}>0.
\end{align*}
But the omitted weak-prior comparison has the opposite sign:
\[
c_H-B_0(1/2)=\frac{47}{10}-\frac{1153}{240}=-\frac5{48}<0.
\]
Orders at $r_0$ are still uniquely zero by (H2).
The posterior is therefore $1/2$, where both cost types prepare.
Entry at $r_0$ is one, and entry at $r_1$ is $\reg{entry_r1_B}$.
Entry falls, so the intended comparison fails. \paper{Appendix~A.3}

For uniform incumbents, all five inequalities in (H1)--(H3) concern continuous functions of the primitives.
Their intersection is open in the domain $0<\beta<1/2$.
The structural restriction $\beta\le1/3$ makes that intersection relatively open within the model domain.
The benchmark uses $\beta<1/3$, so it is also an interior point of that restriction.
Its five positive margins are
\[
\begin{array}{c|c}
c_H-B_0(\pi)&\reg{margin_weakcost_B}\\
B_1(M)-c_H&\reg{margin_recruit_B}\\
k-\Delta_0&\reg{margin_kweak_B}\\
\rho\kappa\Delta_1-k&\reg{margin_kstrong_B}\\
c_H-B_2(M)&\reg{margin_return_B}
\end{array}
\]
These margins prove that the set is nonempty and contains a neighborhood of the benchmark.
The exact slot band also gives a nonempty interval at each admissible $\beta$.
The small margin at the ceiling measures sensitivity, rather than an equality.

\begin{panelnote}[Why no tie rule matters for Proposition 2]
The weak investor comparisons are strict because $\Delta_0<k$.
The two stronger investor comparisons are strict because $k<\rho\kappa\Delta_1\le\rho\kappa\Delta_2$.
All high-cost comparisons at the three equilibrium beliefs are also strict.
Thus neither the investor rule nor the high-cost challenger rule changes the three outcomes.
The fixed zero-cost preparation convention still supplies the entry floor.
Boundary cases outside the strict hypotheses can depend on tie rules, as the previous subsection shows.
\end{panelnote}

Run \texttt{python3 slot/study\_guide/verify/06\_main\_result.py} from the repository root.
It prints the three tuples, auction quantities, entry, ownership, exact band, prior window, ceiling, and H1 counterexample.
